\documentclass[11pt]{article}
\usepackage[a4paper,margin=25mm]{geometry}
\usepackage{amsmath,amssymb,amsthm,hyperref}
\newtheorem{theorem}{Theorem}\newtheorem{lemma}{Lemma}\newtheorem{corollary}{Corollary}
\newcommand{\PP}{\mathcal P}
\title{A polynomial quotient bound across thin density boxes}
\author{Complete proof; independent reviews passed}\date{30 September 2026}
\begin{document}\maketitle
This note makes the interval widths in the positive-completion argument
quantitative. In particular, no condition number of a raw truncated-power
basis is used.

Let $m\ge2$, and let a probability density $\beta$ be supported on $K$
disjoint boxes $J_q=[c_q-\epsilon,c_q+\epsilon]$, each of mass $1/K$.
Assume the complementary gaps, including the two endpoint gaps, have
length at least $\Delta>0$. Assume also the polynomial sampling bound
\begin{equation}\label{sampling}
 K^{-1}\sum_q |p(c_q)|^2\ge c_0\|p\|_{L^2[0,1]}^2
 \qquad(p\in\PP_m),
\end{equation}
where $c_0>0$ is absolute. Perturbed midpoint grids with $K\ge Cm^4$
satisfy this assumption. No bound on $\|\beta\|_\infty$ is needed.
Set
\[
 S_s(a)=\int_{b>a}b^s\beta(b)\,db,\qquad
 V=\PP_m+\operatorname{span}\{a^rS_s(a):r+s\le m\}.
\]
Write $G$ for the union of the complementary gaps, and $\Pi_j$ for
orthogonal projection in $L^2(G)$ onto polynomials of degree at most $j$
separately on each gap. Define
\[
 \mathcal S_0=\operatorname{span}(1,S_0,\ldots,S_m),\qquad
 \mathcal S_1=\operatorname{span}(1,a,S_0,\ldots,S_m,
                                  aS_0,\ldots,aS_{m-1}).
\]

\begin{theorem}\label{quotient}
For sufficiently small absolute $c>0$, if
$\epsilon\le c\Delta/(m+1)^3$, then for every $F\in V$
\begin{align}
 \operatorname{dist}_{L^2[0,1]}(F,\mathcal S_0)
  &\le C(m+1)^4\Delta^{-3/2}\|(I-\Pi_0)F\|_{L^2(G)},\label{q0}\\
 \operatorname{dist}_{L^2[0,1]}(F,\mathcal S_1)
  &\le C(m+1)^6\Delta^{-5/2}\|(I-\Pi_1)F\|_{L^2(G)}.\label{q1}
\end{align}
In particular the kernel of restriction followed by $I-\Pi_j$ is
exactly $\mathcal S_j$, for $j=0,1$.
\end{theorem}
\begin{proof}
Choose any representation $F=P+\sum c_{rs}a^rS_s$. Define $T$ on this
representation by differentiating $P$ and the explicit powers $a^r$
while leaving every $S_s$ unchanged. For $0\le j\le m$ put
\[
 D_j(a)=\sum_{r\ge j,s}c_{rs}(r)_j a^{r-j+s}.
\]
It has degree at most $m-j$, and the absolutely continuous identities
\begin{equation}\label{derivative}
 (T^jF)'=T^{j+1}F-\beta D_j
\end{equation}
hold, with $T^{m+1}F=0$. On the gaps, $T^jF=F^{(j)}$.
Let
\[
 M_j=\|T^jF\|_{L^\infty[0,1]},\qquad
 E_j=\|F^{(j)}\|_{L^\infty(G)}.
\]
Integrating~\eqref{derivative} across a box gives
\begin{equation}\label{boxaverage}
 \left|\int_{J_q}\beta D_j\right|
       \le2E_j+2\epsilon M_{j+1}.
\end{equation}
For any $p\in\PP_m$, Markov and the $L^2$ Nikolskii inequality give
$\|p'\|_\infty\le C(m+1)^3\|p\|_2$. Consequently
\[
 \left|K\int_{J_q}\beta p-p(c_q)\right|
       \le C\epsilon(m+1)^3\|p\|_2.
\]
Assumption~\eqref{sampling} and sufficiently small
$\epsilon(m+1)^3$ imply
\[
 \|p\|_2\le C K\max_q\left|\int_{J_q}\beta p\right|.
\]
Applied to~\eqref{boxaverage}, followed by Nikolskii, this yields
\[
 \|D_j\|_\infty\le C(m+1)K(E_j+\epsilon M_{j+1}).
\]
Integrating~\eqref{derivative} from a box endpoint to an arbitrary
point of that box, and using its total mass $1/K$, now gives
\begin{equation}\label{recurrence}
 M_j\le C(m+1)(E_j+\epsilon M_{j+1}).
\end{equation}
On each gap the usual sup-norm Markov inequality gives
$E_{j+1}\le C(m+1)^2\Delta^{-1}E_j$.
Starting from $M_{m+1}=0$, substitute this inequality in the backward
iteration of~\eqref{recurrence}. Our bound on $\epsilon$ makes the
resulting geometric series summable, so
\begin{equation}\label{globalderivatives}
 M_j\le C(m+1)E_j\qquad(0\le j\le m).
\end{equation}

For~\eqref{q0}, let $U_S$ denote the same linear combination of the
$S_s$ as the monomials in $D_0$. Then $U_S'=-\beta D_0$, so
$(F-U_S)'=TF$. Add a constant to $U_S$ to remove the mean of the
difference. Poincare and~\eqref{globalderivatives} give distance at
most $C(m+1)E_1$. Gapwise Nikolskii and inverse Markov give
\[
 E_1\le C(m+1)\Delta^{-1/2}\|F'\|_{L^2(G)}
 \le C(m+1)^3\Delta^{-3/2}\|(I-\Pi_0)F\|_{L^2(G)}.
\]
This proves~\eqref{q0}.

For~\eqref{q1}, put $U=D_0-aD_1$ and $W=D_1$, and form
$H=U_S+aW_S\in\mathcal S_1$ by replacing powers by the corresponding
$S_s$. Then
\[
 H'=W_S-\beta D_0,\qquad
 H''=-\beta D_1-(\beta D_0)',\qquad
 F''=T^2F-\beta D_1-(\beta D_0)'.
\]
These identities hold distributionally even if $\beta$ jumps at box
endpoints. Thus $(F-H)''=T^2F$. Add a global affine polynomial to $H$
and use the second-order Poincare inequality to obtain distance at
most $CM_2\le C(m+1)E_2$. Finally,
\[
 E_2\le C(m+1)\Delta^{-1/2}\|F''\|_{L^2(G)}
 \le C(m+1)^5\Delta^{-5/2}\|(I-\Pi_1)F\|_{L^2(G)}.
\]
Here subtracting the affine projection on each gap leaves $F''$
unchanged. This proves~\eqref{q1}. The reverse kernel inclusions are
immediate because every $S_s$ is constant on every gap.
\end{proof}

\begin{corollary}[Mass-and-first-moment positive completion]\label{completion}
Assume $\beta$ is rational piecewise polynomial and all endpoints are rational.
Suppose $g$, extended by zero on the boxes, is a rational piecewise
affine function on the gaps, satisfies $1/2\le g\le3/2$ there, and
\[
 \int_0^1 (1-g)H=0\qquad(H\in\mathcal S_1).
\]
If
\[
 C(m+1)^7\Delta^{-3}\|1-g\|_{L^2[0,1]}\le1/4,
\]
then there is a rational piecewise polynomial density $f$, supported
on the gaps, with $1/4\le f\le7/4$ on $G$, having the same mass and first
moment as $g$ on every gap, and satisfying
\[
 \int fF=\int F\qquad(F\in V).
\]
Every polynomial piece has degree at most $m$.
\end{corollary}
\begin{proof}
On $W=(I-\Pi_1)V\subset L^2(G)$ define
$\Lambda((I-\Pi_1)F)=\int(1-g)F$.
The kernel statement and assumed compatibility make this well-defined.
Theorem~\ref{quotient} gives its norm at most
$C(m+1)^6\Delta^{-5/2}\|1-g\|_2$.
Its finite-dimensional Riesz representative $h\in W$ has this same
norm and has zero mass and first moment on each gap. Gapwise
Nikolskii gives
$\|h\|_\infty\le C(m+1)\Delta^{-1/2}\|h\|_2\le1/4$.
Set $f=g+h$. All assertions follow. A rational basis for $W$ has
rational Gram matrix and right side, so the unique Riesz representative
has rational coefficients.
\end{proof}

\begin{theorem}[A fixed number of cell moments]\label{generalq}
Fix $q\ge0$ and suppose $m\ge\max(2,q)$. Retain all box, sampling,
and width assumptions above. Define
\[
 \mathcal S_q=\PP_q+
 \operatorname{span}\{a^rS_s(a):0\le r\le q,\ r+s\le m\},
\]
and let $\Pi_q$ project onto $\PP_q$ separately on each gap. Then
\[
 \operatorname{dist}_{L^2[0,1]}(F,\mathcal S_q)
 \le C_q(m+1)^{2q+4}\Delta^{-q-3/2}
                       \|(I-\Pi_q)F\|_{L^2(G)}\qquad(F\in V).
\]
In particular the kernel is exactly $\mathcal S_q$.
For the completion statement assume additionally that $\beta$ is rational
piecewise polynomial and all endpoints are rational.
If a rational piecewise polynomial $g$ of degree at most $q$ on the
gaps, extended by zero on the boxes, has values between $1/2$ and
$3/2$ and satisfies $\int(1-g)H=0$ for all $H\in\mathcal S_q$,
then the condition
\[
 C_q(m+1)^{2q+5}\Delta^{-q-2}\|1-g\|_{L^2[0,1]}\le1/4
\]
suffices for a rational positive completion $f$, with the same first
$q$ moments as $g$ on every gap and with $\int fF=\int F$ for all
$F\in V$. On each gap $f$ is polynomial of degree at most $m$ and
lies between $1/4$ and $7/4$.
\end{theorem}
\begin{proof}
The already proved backward recurrence and bound
$M_j\le C(m+1)E_j$ hold in all derivative orders through $m$.
For a generator $a^rS_s(a)$, subtract
\[
 H_{r,s}(a)=\int_{b>a}b^s
       \left[\sum_{l=0}^{\min(q,r)}\binom rl b^{r-l}(a-b)^l\right]
       \beta(b)\,db.
\]
Expanding the bracket shows that it belongs to $\mathcal S_q$:
each term is a constant times $a^j S_{r+s-j}(a)$ with $j\le q$.
The bracket is the degree-$q$ Taylor polynomial of $a^r$ about $a=b$.
The residual kernel therefore has zero boundary derivatives through
order $q$. Repeated differentiation under this Volterra integral,
or its distributional version for merely integrable $\beta$, gives
\[
 (a^rS_s-H_{r,s})^{(q+1)}=T^{q+1}(a^rS_s).
\]
Apply the construction to the chosen representation of $F$ and leave
its ordinary polynomial part in the remainder. We obtain
$H\in\mathcal S_q$ with $(F-H)^{(q+1)}=T^{q+1}F$.
Subtract a global degree-$q$ polynomial to remove the first $q+1$
initial Taylor data of $F-H$. The order-$(q+1)$ Poincare inequality
then bounds its $L^2$ norm by $C_q M_{q+1}$.
For $q=m$ the remainder is identically zero; otherwise the recurrence
bounds this by $C_q(m+1)E_{q+1}$.
Gapwise polynomial evaluation and inverse Markov give
\[
 E_{q+1}\le C(m+1)\Delta^{-1/2}\|F^{(q+1)}\|_{L^2(G)}
 \le C_q(m+1)^{2q+3}\Delta^{-q-3/2}
                         \|(I-\Pi_q)F\|_{L^2(G)}.
\]
This proves the quotient and its kernel assertion.

For completion repeat the Riesz argument on
$W=(I-\Pi_q)V$. The functional induced by $\int(1-g)F$ is well-defined
by the kernel assertion and compatibility. Its representative has norm
at most the displayed quotient constant times $\|1-g\|_2$ and zero
moments through degree $q$ in each gap. Gapwise evaluation costs only
$C(m+1)\Delta^{-1/2}$, giving the stated supremum bound. Add this
representative to $g$. Rationality follows from a rational basis,
Gram matrix, and right side, exactly as in Corollary~\ref{completion}.
\end{proof}

\paragraph{How compatibility is checked.}
If $g$ has cumulative mass $r_{0q}$ and cumulative first moment
$r_{1q}$ just before box $q$, then compatibility on the generators
$S_s$ and $aS_s$ is respectively
\[
 \sum_q r_{0q}\int_{J_q}b^s\beta(b)db=\int b^{s+1}\beta(b)db,
 \qquad
 \sum_q r_{1q}\int_{J_q}b^s\beta(b)db=\frac12\int b^{s+2}\beta(b)db.
\]
The remaining two conditions are total mass $1$ and first moment
$1/2$. All of these are linear weighted-moment identities.

\paragraph{Polynomial scales suffice.}
On a gap of length at least $\Delta$, the affine density with prescribed
mass and first moment differs from $1$ by at most $C\Delta^{-2}$ times
the discrepancies of these two moments. Thus if the cumulative data
are within $\eta$ of uniform prefix data and the box centers within
$\zeta$ of the reference nodes, then
\[
 \|g-1\|_{L^\infty(G)}\le C\Delta^{-2}(\eta+\zeta+\epsilon),\quad
 \|1-g\|_{L^2[0,1]}\le
 C\Delta^{-2}(\eta+\zeta+\epsilon)+\sqrt{2K\epsilon}.
\]
When $K,\Delta^{-1}$ are polynomial in $m$, choosing successively small
fixed inverse powers for $\eta,\epsilon,\zeta$ meets all hypotheses.
This yields polynomially small boxes, rather than an unspecified
continuity scale.
\end{document}
